\chapter{Background and State of the Art}
\label{chap:background}

The rapid proliferation of Low Earth Orbit (LEO) satellite constellations has fundamentally redefined the landscape of satellite communications and space-borne computing. Transitioning from traditional passive data relays to active processing nodes requires a comprehensive understanding of orbital dynamics, task offloading models, and orchestration strategies. This chapter provides the theoretical foundation and contextual literature for this thesis. It details the paradigms of Orbital Edge Computing (OEC), analyzes task offloading and routing dynamics in LEO networks, compares centralized and distributed orchestration models, reviews state-of-the-art baseline solutions, and articulates the research gap addressed by our work.

\section{Satellite Computing Paradigms}
\label{sec:bg_paradigms}

Space-based networking architectures have historically been categorized by their orbital altitudes: Geostationary Earth Orbit (GEO), Medium Earth Orbit (MEO), and Low Earth Orbit (LEO). GEO satellites, positioned at approximately $35,786 \text{ km}$, offer stationary coverage over vast geographic areas but suffer from high round-trip propagation delays (exceeding $250 \text{ ms}$) and high free-space path loss. MEO satellites ($2,000 \text{ km} - 35,786 \text{ km}$) reduce latency slightly but remain unsuitable for ultra-low-latency processing.

In contrast, LEO satellites operate at altitudes ranging from $300 \text{ km}$ to $2,000 \text{ km}$. Their proximity to Earth drastically lowers propagation latency ($10 - 40 \text{ ms}$ round-trip) and reduces power loss, making them ideal for high-throughput, real-time communications. However, LEO satellites travel at extreme velocities relative to the surface ($\sim 7.5 \text{ km/s}$), completing an orbit in approximately $90 - 120 \text{ minutes}$. This rapid mobility results in dynamic topologies where individual satellites remain in line-of-sight with a fixed ground position for only a few minutes.

To process data efficiently within this dynamic environment, satellite network architectures have evolved through three primary paradigms:

\begin{enumerate}
    \item \textbf{Transparent Bent-Pipe Relay:} Satellites act merely as transparent repeaters. Raw data captured by Earth Observation (EO) sensors or transmitted by Ground Users (GUs) is amplified and forwarded directly to Ground Stations (GSs). Computation occurs entirely in terrestrial cloud centers. This model suffers from severe downlink bottlenecks and complete connectivity loss when flying over remote regions (e.g., oceans or polar zones) lacking GS infrastructure.
    \item \textbf{Terrestrial Cloud Offloading with Inter-Satellite Routing:} Satellites utilize Inter-Satellite Links (ISLs) to route raw data across orbital planes until a satellite with an active GS connection is reached. While this mitigates local GS coverage gaps, the volume of raw data still strains optical and RF link bandwidth.
    \item \textbf{Orbital Edge Computing (OEC) / Satellite Edge Computing (SEC):} Satellites are equipped with onboard processing units—termed Satellite Edge Nodes (SENs)—capable of executing compute-intensive algorithm payloads (e.g., AI inference, image compression, signal filtering) directly in orbit. Raw data is processed locally at the source node or offloaded across ISLs to adjacent SENs. Only refined, actionable insights or small computation results are downlinked to Earth, drastically reducing bandwidth consumption and end-to-end response times.
\end{enumerate}

\section{Task Offloading and Routing in LEO Networks}
\label{sec:bg_offloading_routing}

Task offloading in Orbital Edge Computing involves receiving task execution requests generated by Ground Users (GUs) or space-borne sensors and strategically allocating their computational execution across available SENs while routing payloads through Inter-Satellite Links (ISLs).

\subsection{System Entities and Network Topology}

An OEC network architecture consists of three core structural entities:
\begin{itemize}
    \item \textbf{Ground Users (GUs) / Remote Sensors:} Distributed terrestrial entities that issue computational task requests $r_i$. Each request is characterized by a data payload size $D^{data}_r$ (in bits), required CPU compute cycles $W_r$, and a strict operational deadline $D_r$ (in seconds).
    \item \textbf{Satellite Edge Nodes (SENs):} LEO satellites equipped with CPU cores, memory modules, and onboard energy reserves (solar panels and battery banks). Each SEN $s_j$ possesses a processing capacity $f_j$ (cycles/second) and an operational energy budget $E_j$.
    \item \textbf{Inter-Satellite Links (ISLs):} High-speed communication channels connecting adjacent satellites. Intra-plane ISLs connect satellites within the same orbital plane (stable distance), whereas inter-plane ISLs connect satellites in adjacent planes (dynamically varying distance and pointing angles). Each ISL link $(u, v)$ is constrained by bandwidth $B_{uv}$ and propagation delay $\tau_{uv}$.
\end{itemize}

\subsection{Workload Profiles and Constraints}

Workloads submitted to an OEC network are typically heterogeneous, categorized into two main profiles:
\begin{itemize}
    \item \textbf{Real-Time Workloads:} Tasks generated by time-critical applications (e.g., emergency disaster detection, defense alerts) requiring immediate processing with tight deadlines $D_r$.
    \item \textbf{Batch / Delay-Tolerant Workloads:} Tasks associated with routine telemetry, scientific processing, or background data aggregation with relaxed deadlines.
\end{itemize}

Offloading decisions must jointly satisfy physical compute constraints ($\sum W_r \le f_j$), network bandwidth limits, energy consumption limits (computation power $P_{comp}$ and transmission power $P_{tx}$), and strict end-to-end deadline bounds ($T_{total} = T_{propagation} + T_{transmission} + T_{queue} + T_{computation} \le D_r$).

\section{Centralized vs. Distributed Orchestration}
\label{sec:bg_centralized_vs_distributed}

Managing resource allocation across hundreds of moving SENs requires an orchestration paradigm. Literature categorizes these paradigms into two primary architectures: Centralized and Distributed.

\subsection{Centralized Orchestration}

In a centralized orchestration architecture, a dedicated master entity—such as a designated lead satellite (Ground/Orbital Orchestrator) or a ground control station—collects global network state information (CPU load, battery state, link traffic) from all SENs.

\begin{itemize}
    \item \textbf{Advantages:} The central orchestrator maintains complete global visibility over all network resources, enabling mathematically optimal decision-making. Global optimization models (e.g., Integer Linear Programming) can eliminate load imbalances, prevent local hotspots, and minimize network-wide energy consumption.
    \item \textbf{Disadvantages:} Centralized schemes incur significant control signaling overhead to aggregate state information. Furthermore, requests must travel to the orchestrator, causing queuing delays and introducing a single point of failure and bottleneck under high request arrival rates ($\lambda$).
\end{itemize}

\subsection{Distributed Orchestration}

In distributed orchestration, individual SENs or localized clusters make autonomous offloading and routing decisions based solely on locally observed state information or partial neighborhood telemetry.

\begin{itemize}
    \item \textbf{Advantages:} Distributed schemes eliminate centralized control bottlenecks, offer sub-millisecond decision latency, and demonstrate high fault tolerance against dynamic topology changes and satellite failures.
    \item \textbf{Disadvantages:} Lacking global visibility, distributed algorithms often converge to sub-optimal solutions. They are prone to localized network congestion, unbalanced CPU utilization, and premature battery depletion on popular forwarding paths.
\end{itemize}

\section{State-of-the-Art Approaches}
\label{sec:bg_state_of_the_art}

To place our research in context, we review key baseline algorithms and mathematical models established in recent IEEE literature on Orbital Edge Computing and Task Offloading.

\subsection{Distributed Heuristic Frameworks}

Heuristic approaches sacrifice theoretical optimality in favor of computational speed and low decision latency:

\begin{itemize}
    \item \textbf{OrbitAware Architecture:} A dynamic heuristic framework designed specifically for highly dynamic LEO topologies. \textit{OrbitAware} leverages periodic orbital trajectory updates to predict link availability and select shortest processing/routing paths locally, reducing control overhead at the cost of potential link queue saturation.
    \item \textbf{DTS-base (Dynamic Task Scheduling - Base):} A greedy distributed scheduling baseline that assigns incoming tasks to the nearest available SEN with sufficient idle CPU capacity. While fast, it ignores global energy budgets and multi-hop link traffic.
    \item \textbf{DTS-APopt (Dynamic Task Scheduling - Access Point Optimized):} An enhanced version of \textit{DTS-base} that incorporates multi-criteria decision metrics, optimizing both the local access point selection and initial ISL hop bandwidth before dispatching tasks.
\end{itemize}

\subsection{Optimization-Based and Distributed ILP Formulations}

To improve allocation quality, recent studies have formulated mathematical optimization models using Integer Linear Programming (ILP) and mixed-integer techniques:

\begin{itemize}
    \item \textbf{Distributed / Hierarchical ILP:} Frameworks that decompose a global ILP problem into smaller sub-problems solved locally across satellite clusters. While more scalable than pure centralized ILP, distributed ILP approaches suffer from convergence delays and coordination overhead during inter-cluster task transfers.
    \item \textbf{Weighted Objective Functions:} Optimization models in the literature frequently balance competing goals—such as latency minimization versus energy preservation—by assigning scalar weights ($W_{lat}, W_{en}$) to objective terms.
\end{itemize}

\section{Research Gap and Thesis Positioning}
\label{sec:bg_research_gap}

While state-of-the-art distributed heuristics offer rapid processing times and distributed ILP frameworks attempt local decomposition, a significant research gap remains regarding the practical limits of \textbf{Centralized ILP Orchestration} in LEO satellite networks.

Existing literature frequently assumes either ideal zero-latency centralized orchestrators or dismisses centralized optimization entirely due to its computational complexity. Few works rigorously examine the trade-offs when a centralized ILP orchestrator is equipped with an integrated \textbf{Dynamic Batching} and \textbf{Timeout Mechanism} ($BS / BT$) to manage request queues under heavy workloads ($\lambda > 6 \text{ req/s}$).

This thesis addresses this exact gap by:
\begin{enumerate}
    \item Formalizing a centralized ILP model that dynamically toggles between latency minimization (\textit{Mapping Time}) and energy conservation (\textit{Mapping Energy}).
    \item Quantifying the systemic performance bounds of centralized optimization when coupled with dynamic request batching and queuing delay compensation.
    \item Conducting a rigorous comparative benchmark against established state-of-the-art heuristics (\textit{OrbitAware}, \textit{DTS-base}, \textit{DTS-APopt}) and distributed ILP variants under realistic discrete-event simulation conditions (\texttt{SimPy}).
\end{enumerate}